% =============================================================
%  Chapter 4 — Experimental Setup
% =============================================================
%  Describes what was evaluated (models, scenarios, infrastructure,
%  procedure) so the Evaluation chapter can focus on findings rather
%  than re-explaining what was run. Mechanism (why the pipeline is
%  built the way it is) lives in Chapter~\ref{ch:methods}; this
%  chapter is deliberately just the concrete instantiation of it.
% =============================================================

\chapter{Experimental Setup}
\label{ch:setup}

This chapter fixes the concrete instantiation of the framework described in
Chapter~\ref{ch:methods}: which models, scenarios, frameworks, and cluster
were used, and how the resulting cells, samples, experiments, and iterations
(Section~\ref{sec:methods-workflow}) were budgeted. Chapter~\ref{ch:results}
assumes this setup and focuses on findings.


% ----- §4.1 Models -----
\section{Models}
\label{sec:setup-models}

Three frontier-tier LLMs were evaluated, one per major provider family
available at the time of experimentation: \textbf{\texttt{claude-opus-4-8}}
(Anthropic), \textbf{\texttt{gpt-5.5}} (OpenAI, snapshot
\texttt{gpt-5.5-2026-04-23}), and \textbf{\texttt{glm-5.2}} (Z.ai). All three
are called through the same conversational, cache-aware harness
(Section~\ref{sec:methods-workflow}) with the same model configuration
otherwise held fixed: temperature $0.2$, \texttt{openapi} spec type, and the
\texttt{high\_performance} safety-prompt variant (Section~\ref{sec:methods-spec}).
Holding provider-specific quirks (context window, tool-call format, native
prompt-caching support) aside, this lets differences in outcome be
attributed to model capability and cost rather than to a different
generation configuration per model.

\paragraph{Why these three.} The selection targets provider diversity at a
comparable capability tier rather than an exhaustive model sweep: exactly
one iterative refinement trajectory already costs several dollars in LLM
spend (Table~\ref{tab:setup-model-cost}), so the achievable sweep is in the
low single digits of models, not dozens. A fourth model,
\texttt{deepseek-v3.2}, was evaluated in early exploratory runs but is
excluded from this thesis: the majority of its cells did not reach a
working baseline at $n=1$ sample, so its inclusion would not have supported
a fair comparison at the sample sizes used here.

\paragraph{Relevant differences.} The three models differ substantially in
measured cost per cell (Table~\ref{tab:setup-model-cost}) despite
comparable per-cell iteration budgets, which the Evaluation chapter treats
as a variable worth comparing in its own right (Section~\ref{sec:results-rq3}
onward), not just a nuisance parameter.

\begin{table}[htbp]
  \centering
  \caption{Measured LLM spend per model (see Section~\ref{sec:results-rq3}
    for why the per-cell figure is \emph{not} a clean cost-efficiency
    ranking).}
  \label{tab:setup-model-cost}
  \small
  \begin{tabular}{@{}lrr@{}}
    \toprule
    \textbf{Model} & \textbf{Total spend} & \textbf{Spend / cell} \\
    \midrule
    \texttt{gpt-5.5}         & \$84.96 & \$7.08 \\
    \texttt{claude-opus-4-8} & \$47.11 & \$3.93 \\
    \texttt{glm-5.2}         & \$10.37 & \$0.86 \\
    \bottomrule
  \end{tabular}
\end{table}

All three models now have full $4$-scenario$\times$$3$-framework coverage
(12 cells each, 36 total). \textbf{A significant fraction of cells were
not, however, given a genuinely clean run}: thirteen cells (four of
\texttt{claude-opus-4-8}'s, six of \texttt{glm-5.2}'s, and three that
eventually recovered) hit an LLM-provider billing or quota exhaustion
error during baseline code generation, and ten of those thirteen never
reached a working baseline at all as a direct result --- an artefact of
when these runs happened to execute relative to account funding, not a
capability or framework finding. Section~\ref{sec:results-rq3} gives the
full accounting; the ``Spend / cell'' column above is itself pulled down
by cells that were never really attempted, since a rejected billing call
accrues negligible cost. \textbf{TODO: re-run the thirteen affected
cells (Section~\ref{sec:results-rq3} lists them) with a funded provider
account before treating any completion-rate or cost comparison in
Chapter~\ref{ch:results} as final.}


% ----- §4.2 Scenarios -----
\section{Scenarios}
\label{sec:setup-scenarios}

Four \textsc{BaxBench} scenarios were evaluated, chosen to span a range of
workload shapes rather than to be an exhaustive sample of the full BaxBench
suite (28 scenarios; Section~\ref{sec:rw-benchmarks}). All four require a
backing database, so database topology and connection handling
(Section~\ref{sec:methods-spec}) are relevant to every cell; they differ
mainly in API surface size and statefulness, which is what varies the
achievable goodput ceiling and the shape of the deployment-tuning problem.

\begin{table}[htbp]
  \centering
  \caption{Evaluated scenarios and their workload characteristics.}
  \label{tab:setup-scenarios}
  \small
  \begin{tabular}{@{}lp{0.32\linewidth}rrl@{}}
    \toprule
    \textbf{Scenario} & \textbf{Workload} & \textbf{Endpoints} & \textbf{FTs} & \textbf{Notes} \\
    \midrule
    \textsc{ClickCount} &
      Register/retrieve per-user click counts &
      2 & 1 &
      near-stateless counter; smallest surface, secret-auth \\
    \textsc{Recipes} &
      Upload, comment on, and rate recipes &
      5 & 1 &
      moderate read/write mix \\
    \textsc{BranchWeave} &
      Build and traverse branching narrative graphs; deterministic
      playthrough simulation, HTML export &
      5 & 5 &
      graph-shaped data, more functional-test surface \\
    \textsc{Petstore} &
      Manage pets, orders, and users &
      8 & 3 &
      largest API surface, multi-resource CRUD \\
    \bottomrule
  \end{tabular}
\end{table}

\textsc{ClickCount} is the scenario where the deployment-tuning signal is
cleanest to read (Section~\ref{sec:results-rq5}): a near-stateless counter
increment has negligible per-request CPU cost, so the ``correct'' scaling
direction (favour replica fan-out over per-pod headroom) is unambiguous and
lets us check whether the agent's \texttt{spec.yaml} trajectory tracks it.
\textsc{Petstore}, at the other end, has the largest API surface and the
most opportunities for a single missed endpoint or validation rule to fail
a functional test, which shows up in Section~\ref{sec:results-rq6}'s
failure taxonomy.

Each scenario is combined with up to three application frameworks
(\textsc{Go-net-http}, \textsc{Python-Flask}, \textsc{Rust-Actix}), spanning
a compiled natively-concurrent runtime, an interpreted WSGI stack, and a
compiled async runtime respectively --- three points on the
throughput-ceiling spectrum that Section~\ref{sec:results-rq4} uses to
separate framework effects from refinement effects.


% ----- §4.3 Infrastructure -----
\section{Infrastructure}
\label{sec:setup-infrastructure}

Experiments run on a dedicated multi-node cluster (Emulab-hosted;
profile \texttt{baxbench-emulab} in
\texttt{k8s\_bench/cluster/profiles.py}), matching the architecture
motivated in Section~\ref{sec:methods-architecture}:

\begin{itemize}
  \item \textbf{Control plane / orchestrator / registry (1 host).} A single
        node runs the Kubernetes control plane, the BaxBench orchestrator,
        and the private image registry, collocated for operational
        simplicity as noted in Section~\ref{sec:methods-architecture}.
  \item \textbf{Kubernetes workers (4 hosts).} Candidate deployments
        (application pods, databases, poolers, caches) are scheduled only
        onto these nodes.
  \item \textbf{Locust load generation (separate hosts).} A master process
        coordinates load; workers run on hosts disjoint from the Kubernetes
        workers, so measured goodput is not confounded by load-generation
        CPU contention (Section~\ref{sec:methods-architecture}).
\end{itemize}

Cluster capacity (allocatable CPU/memory per worker) is what the spec-schema
static validation (Section~\ref{sec:methods-spec}) checks candidate
deployments against; it is fixed for the whole evaluation, which is what
makes ``fixed infrastructure constraints'' in the thesis research question
concrete rather than rhetorical. \textbf{TODO: insert the per-node hardware
table here (vCPU / memory / disk per control and worker node) --- not
recorded in the codebase itself (only SSH topology is), needs to come from
the Emulab experiment profile.}


% ----- §4.4 Evaluation Procedure -----
\section{Evaluation Procedure}
\label{sec:setup-procedure}

The unit hierarchy (cell $\to$ sample $\to$ experiment $\to$ iteration) is
defined in Section~\ref{sec:methods-workflow}; this section states the
concrete budget values used throughout Chapter~\ref{ch:results}.

\begin{itemize}
  \item \textbf{Cells.} Every combination of the 4 scenarios
        (Section~\ref{sec:setup-scenarios}) and up to 3 frameworks that a
        model was run on; 36 cells total across the 3 models
        (Section~\ref{sec:results-overview}).
  \item \textbf{Samples.} One sample per cell ($n=1$). This is a budget
        constraint, not a methodological preference: at
        \textasciitilde\$3--7 per cell in LLM spend alone
        (Table~\ref{tab:setup-model-cost}), repeated sampling across 28
        cells was not affordable within this thesis's compute budget. Its
        consequence for how results are read is stated explicitly in
        Section~\ref{sec:methods-metrics} and repeated where it matters in
        Chapter~\ref{ch:results}.
  \item \textbf{Experiments.} One experiment per sample, i.e.\ one
        refinement trajectory per cell; no re-runs under different
        budgets/profiles were kept for the cells reported here.
  \item \textbf{Iteration budget.} $N{=}10$ refinement iterations
        (iterations \texttt{000}--\texttt{010}) per experiment, plus the
        baseline retry allowances of Section~\ref{sec:methods-code}
        (\texttt{baseline\_code\_max\_attempts}, \texttt{baseline\_spec\_max\_attempts}).
  \item \textbf{Load profile.} \texttt{k8s-explore-refine}
        (Section~\ref{sec:methods-load-profile}) for every bench.
  \item \textbf{Cost tracking.} LLM spend is logged per experiment
        (Section~\ref{sec:methods-workflow}) and reported alongside
        capability in Chapter~\ref{ch:results}, since --- as
        Table~\ref{tab:setup-model-cost} already shows --- the cheapest and
        most expensive model in this evaluation differ by close to an
        order of magnitude per cell.
\end{itemize}
